% ==============================================================================
% MÓDULO: CAPÍTULO IV - METODOLOGÍA: DISEÑO FENOMENOLÓGICO DE INVESTIGACIÓN
% Archivo: subcapitulos/sec_04_metodologia.tex
% ==============================================================================
\section{METODOLOGÍA: DISEÑO FENOMENOLÓGICO DE INVESTIGACIÓN}
\label{sec:cap4_metodologia}

\subsection{Fundamentación filosófica y epistemológica del diseño fenomenológico}
\label{sec:4.1_enfoque}
La presente investigación se inserta en el paradigma interpretativo-naturalista y adopta un \textbf{enfoque puramente cualitativo con diseño fenomenológico} (Hernández-Sampieri, Fernández-Collado y Baptista-Lucio, 2014, Cap. 12 y 15) \cite{sampieri2018}. Como afirman Roberto Hernández-Sampieri y colaboradores, la fenomenología no es únicamente un conjunto de técnicas procedimentales, sino simultáneamente una filosofía, una postura epistemológica y un riguroso diseño investigativo cuyo origen se remonta al pensamiento de Edmund Husserl (1859--1938) \cite{husserl1949}. Su propósito cardinal radica en \textbf{explorar, describir y comprender las experiencias de las personas con respecto a un determinado fenómeno y descubrir los elementos en común (la esencia compartida) de tales vivencias} \cite{sampieri2018}.

En el campo de las ciencias de la salud y el cuidado de enfermería, Norlyk y Harder \cite{norlyk2010} evaluaron 88 estudios empíricos demostrando que la fenomenología es el abordaje con mayor pertinencia ontológica para develar las vivencias subjetivas directas frente a la atención asistencial. El estudio se rige por las siguientes distinciones y premisas:

\begin{itemize}
    \item \textbf{Diferenciación frente a la Teoría Fundamentada:} Mientras que la teoría fundamentada persigue generar un modelo teórico explicativo abstracto a partir de los datos, en el diseño fenomenológico \textbf{los investigadores trabajan directamente con las unidades o declaraciones de los participantes y sus vivencias}, explorando lo que los individuos tienen en común respecto al fenómeno, más que abstraerlas para crear un modelo hipotético impuesto \cite{sampieri2018}. Los datos provienen de sentimientos, emociones, razonamientos, percepciones y sentires somáticos (dolor, angustia, consuelo, gratitud, frialdad) \cite{benner2008, sampieri2018}.
    \item \textbf{Diferenciación frente al Diseño Narrativo:} En tanto que el diseño narrativo se focaliza en la sucesión cronológica de eventos (la historia biográfica en el tiempo), el diseño fenomenológico se enfoca con exclusividad en la \textbf{esencia de la experiencia compartida} (lo que vivenciaron los participantes y de qué forma lo vivieron en su corporalidad y contexto) \cite{sampieri2018}.
    \item \textbf{Diferenciación frente a la Etnografía y la Investigación-Acción:} No busca el estudio holístico y prolongado de una cultura comunal completa \cite{cotan2020}, ni introduce intervenciones experimentales o modificaciones deliberadas de la realidad (investigación-acción), preservando la fidelidad descriptiva e interpretativa del fenómeno tal como acontece en su estado natural.
    \item \textbf{Articulación de dos vertientes fenomenológicas complementarias (Sampieri et al., 2014):}
    \begin{enumerate}
        \item \textit{Fenomenología empírica o psicológica \cite{moustakas1994, creswell2013, wilson2007}:} Centrada en describir la esencia de las vivencias de los usuarios; la investigadora suspende temporalmente sus juicios previos (\textit{epojé} o \textit{bracketing}) para aprender de forma pura la experiencia del paciente.
        \item \textit{Fenomenología hermenéutica \cite{vanmanen1990, creswell2013}:} Orientada a la interpretación reflexiva de los discursos y de los ``textos de la vida'', develando significados latentes, silencios y metáforas en quechua chanka dentro del contexto sociocultural andino de Huamanga.
    \end{enumerate}
\end{itemize}

\subsection{Los cuatro existenciales de la experiencia humana vivida (van Manen y Sampieri)}
\label{sec:4.2_existenciales}
Siguiendo las directrices de Roberto Hernández-Sampieri et al. \cite{sampieri2018} (p. 494) y Max van Manen \cite{vanmanen1990}, toda indagación fenomenológica exige contextualizar las experiencias humanas en torno a cuatro estructuras existenciales universales del mundo de la vida (\textit{Lebenswelt}):

\begin{enumerate}
    \item \textbf{Temporalidad (el tiempo vivido):} Alude a la percepción subjetiva del paso del tiempo y al momento vivencial en que acontecen los sucesos. En el Centro de Salud Belén, la temporalidad se vivencia en la vigilia del madrugar a las 4:00 AM, la incertidumbre angustiosa durante las horas en la fila exterior antes de abrir ventanilla, el tiempo detenido en las bancas de espera y la sensación de sosiego o prisa en la consulta.
    \item \textbf{Espacialidad (el espacio vivido):} Comprende el entorno físico, ambiental y simbólico donde transcurre la experiencia. Abarca la intemperie gélida de la acera matutina, el hacinamiento en pasillos angostos, la incomodidad de bancas duras y la atmósfera de intimidad, calidez o invasión a la privacidad dentro del consultorio.
    \item \textbf{Corporalidad (el cuerpo vivido):} Examina la dimensión física y somática de la persona que vive el fenómeno. En los usuarios de Belén, el cuerpo vivido se manifiesta a través del frío penetrante en la madrugada, el cansancio muscular por permanecer horas de pie, la fatiga del ayuno, el llanto del niño en brazos de la madre, el dolor somático que motiva la búsqueda de ayuda y el alivio físico al recibir la atención.
    \item \textbf{Relacionalidad (la relación humana vivida e intersubjetividad):} El encuentro intersubjetivo que se establece entre el profesional de salud y el usuario doliente. Constituye el corazón del cuidado: el saludo afectuoso (\textit{napaykuy}), el contacto visual sincero, la escucha atenta, la calidez o desdén burocrático, y el puente comunicativo forjado mediante el uso fluido de la lengua originaria quechua chanka.
\end{enumerate}

\subsection{Postura epistemológica, reflexividad y reducción fenomenológica (\textit{Epojé})}
\label{sec:4.3_postura}
En la investigación cualitativa fenomenológica, el investigador reconoce que \textbf{su propia persona actúa como el instrumento principal de indagación y contacto humano} (Hernández-Sampieri et al., 2014, Cap. 14) \cite{sampieri2018}. La investigadora adopta una postura reflexiva continua gobernada por:

\begin{enumerate}
    \item \textbf{Práctica de la Reducción Fenomenológica (\textit{Epojé} / \textit{Bracketing}):} De acuerdo con Edmund Husserl \cite{husserl1949}, Clark Moustakas \cite{moustakas1994} y John Creswell \cite{creswell2013}, la investigadora realiza una suspensión consciente y sistemática de sus propios juicios previos, experiencias previas como interna de enfermería y conceptos clínicos anticipados, para aproximarse a los relatos de los usuarios con total apertura y asombro epistémico.
    \item \textbf{Construcción de Clima de Confianza (\textit{Rapport}):} Establecimiento de una relación horizontal y fraterna con los participantes. En los entornos sanitarios andinos, el informante sólo comparte sus temores profundos y sus percepciones íntimas si experimenta un trato genuinamente afectuoso y desprovisto de jerarquías institucionales.
    \item \textbf{Empatía Intercultural y Sensibilidad Lingüística:} Comprensión del mundo afectivo del usuario andino desde su propia cosmovisión, interactuando en quechua chanka o castellano para que el participante se exprese con libertad, sin temor a ser juzgado o desatendido en su salud.
\end{enumerate}

\subsection{Muestreo cualitativo: Muestra de casos tipo, homogénea y saturación teórica}
\label{sec:4.4_muestra}
En consonancia con el Capítulo 13 de Roberto Hernández-Sampieri, Carlos Fernández-Collado y Pilar Baptista-Lucio \cite{sampieri2018}, el muestreo cualitativo no persigue la representatividad matemática de una población ni el cálculo de tamaños muestrales mediante fórmulas probabilísticas; su propósito fundamental es \textbf{seleccionar casos que aporten profundidad, riqueza informativa y comprensión exhaustiva del fenómeno}:

\begin{itemize}
    \item \textbf{Tipo de Muestra Dirigida:} Se adopta una \textbf{muestra de casos tipo y homogénea}, constituida por personas usuarias que poseen perfiles comunes y que han vivenciado de manera directa, prolongada y personal el fenómeno del acceso y la atención en el Centro de Salud Belén.
    \item \textbf{Principio de Saturación de Categorías:} El número definitivo de participantes no se fija de manera rígida apriorística, sino que se define en el trabajo de campo por el principio de \textbf{saturación teórica} \cite{esbensen2008}, entendido como el momento metodológico en el cual la recolección de nuevas entrevistas ya no genera propiedades, comprensiones ni categorías emergentes novedosas sobre la vivencia asistencial. Se estima preliminarmente una muestra referencial de 12 a 18 participantes informantes clave.
\end{itemize}

\textbf{Criterios de Inclusión Fenomenológicos:}
\begin{enumerate}
    \item Personas usuarias mayores de 18 años, de ambos sexos, que hayan recibido atención en los consultorios de Medicina General, Enfermería (CRED, Inmunizaciones), Obstetricia u Odontología del Centro de Salud Belén en los últimos tres meses.
    \item Personas que hayan vivenciado personalmente la trayectoria de acceso (madrugar, espera en salas y contacto clínico).
    \item Participación voluntaria ratificada mediante la firma o huella dactilar del Consentimiento Informado (\hyperref[anx:4_consentimiento]{Anexo 4}).
    \item Capacidad de verbalizar y reflexionar sobre sus vivencias asistenciales en quechua chanka o castellano.
\end{enumerate}

\textbf{Criterios de Exclusión Fenomenológicos:}
\begin{enumerate}
    \item Personas con afecciones cognitivas agudas o alteraciones de la comunicación verbal que imposibiliten el relato dialógico reflexivo.
    \item Usuarios que acudan únicamente a gestiones administrativas breves de ventanilla sin haber recibido atención clínica en consultorio.
\end{enumerate}

\subsection{Técnicas e instrumentos de recolección de vivencias (Sampieri, Cap. 14)}
\label{sec:4.5_tecnicas}
Definiendo al \textbf{investigador como el instrumento principal de recolección de los datos}, se aplican técnicas cualitativas no estandarizadas que permiten capturar el fenómeno en su ambiente natural:

\begin{enumerate}
    \item \textbf{Entrevistas Fenomenológicas en Profundidad:} Técnica central (sesiones de 45 a 60 minutos) conducida de forma abierta y semiestructurada a través de una guía flexible (\hyperref[anx:1_guia]{Anexo 1}), orientada a profundizar en el significado, la corporalidad, el tiempo y los lazos interpersonales experimentados en el centro de salud \cite{creswell2013, norlyk2010}.
    \item \textbf{Grupos de Enfoque (Focus Groups):} Encuentros colectivos dialógicos (6 a 8 usuarios por grupo) en un ambiente tranquilo cedido por el centro de salud, diseñados para confrontar, enriquecer y construir colectivamente las categorías vivenciales compartidas respecto a las filas de madrugada y el trato humano.
    \item \textbf{Observación Cualitativa no Estructurada e Inmersión en el Ambiente Natural:} Presencia atenta y reflexiva de la investigadora en los escenarios naturales (la acera exterior a las 4:30 AM, los pasillos de espera, las ventanillas y el ingreso a consultorios), prestando atención a expresiones espontáneas, posturas de fatiga y dinámicas de interacción.
    \item \textbf{Bitácora y Diario de Campo Reflexivo:} Instrumento fundamental donde la investigadora registra notas de campo, impresiones sensoriales, silencios, expresiones corporales, inflexiones de voz en quechua y notas de reflexividad pos-entrevista.
    \item \textbf{Triangulación Múltiple (Métodos y Fuentes de Datos):} Conforme a Sampieri et al. \cite{sampieri2018} y Okuda y Gómez-Restrepo \cite{okuda2005}, se efectúa una triangulación de métodos (entrevistas en profundidad + grupos de enfoque + observación no estructurada + bitácora de campo) y una triangulación de fuentes de datos mediante tres miradas concurrentes:
    \begin{itemize}
        \item \textit{Usuarios pacientes dolientes:} Madres en CRED, gestantes y adultos mayores (la vivencia sentida en carne propia).
        \item \textit{Personal asistencial y de enfermería:} Quienes proveen el cuidado (retos operativos y condicionantes institucionales).
        \item \textit{Acompañantes y líderes comunitarios de Belén:} Familiares y vecinos que presencian las esperas (la mirada comunitaria colectiva).
    \end{itemize}
\end{enumerate}

\subsection{Criterios de rigor científico cualitativo (Hernández-Sampieri, Cap. 14)}
\label{sec:4.6_rigor}
En la investigación cualitativa se reemplazan radicalmente los conceptos positivistas de validez y confiabilidad estadística por los cuatro criterios de rigor cualitativo establecidos por Roberto Hernández-Sampieri, Carlos Fernández-Collado y Pilar Baptista-Lucio \cite{sampieri2018} (articulando los postulados clásicos de Lincoln y Guba \cite{gubalincoln1985}):

\begin{table}[H]
\centering
\small
\begin{tabularx}{\textwidth}{p{3.2cm}Xp{5.5cm}}
\toprule
\textbf{Criterio de Rigor} & \textbf{Definición Metodológica (Sampieri)} & \textbf{Estrategia Operativa de Aseguramiento} \\
\midrule
\textbf{1. Dependencia (Consistencia Cualitativa)} & Coherencia interna, estabilidad y consistencia en los procedimientos de recolección y sistematización de los datos. & $\bullet$ Pistas de auditoría transparentes.\newline $\bullet$ Registro detallado en la bitácora de campo.\newline $\bullet$ Protocolo metodológico sistemático de transcripción y codificación. \\
\addlinespace
\textbf{2. Credibilidad (Validez Interna Cualitativa)} & Fidelidad y veracidad con la que las interpretaciones y narrativas del investigador capturan los significados reales de los participantes. & $\bullet$ Triangulación múltiple de métodos y fuentes.\newline $\bullet$ Transcripción íntegra fonética (\textit{verbatim}).\newline $\bullet$ Verificación con los participantes (\textit{member checking} / confirmación testimonial). \\
\addlinespace
\textbf{3. Transferencia (Aplicabilidad Contextual)} & Posibilidad de transferir o contextualizar los hallazgos a otros escenarios con características ecológicas y sociales análogas. & $\bullet$ Descripción densa y detallada del contexto del C.S. Belén (Capítulo II).\newline $\bullet$ Caracterización sociocultural y lingüística minuciosa de los informantes. \\
\addlinespace
\textbf{4. Confirmabilidad (Neutralidad Reflexiva)} & Garantía de que los hallazgos representan genuinamente las vivencias de los informantes y no los sesgos, prejuicios o intereses del investigador. & $\bullet$ Práctica sistemática de la reducción fenomenológica (\textit{epojé} o \textit{bracketing}).\newline $\bullet$ Memos analíticos de reflexividad.\newline $\bullet$ Auditoría y revisión crítica con el asesor Dr. Manglio Aguirre Andrade. \\
\bottomrule
\end{tabularx}
\caption{Criterios de rigor científico cualitativo según Hernández-Sampieri (2014).}
\label{tab:rigor_cualitativo}
\end{table}

\subsection{Procedimiento sistemático de análisis cualitativo fenomenológico (Sampieri, Cap. 14 y 15)}
\label{sec:4.7_analisis}
El análisis cualitativo se regirá por la ruta procedimental del diseño fenomenológico establecida por Roberto Hernández-Sampieri y colaboradores \cite{sampieri2018} (Capítulos 14 y 15), articulada con las directrices canónicas de Colaizzi, Moustakas \cite{moustakas1994} y Creswell \cite{creswell2013}, apoyada técnicamente mediante el software cualitativo ATLAS.ti v9:

\begin{enumerate}
    \item \textbf{Organización, transcripción y lectura holística:}
    Preparación y digitalización exhaustiva de los registros sonoros en transcripciones literales palabra por palabra (\textit{verbatim}) dentro de las 24 horas posteriores a cada sesión, preservando pausas, llanto, silencios y vocablos en quechua chanka. A continuación, inmersión y lecturas repetidas del material para alcanzar una aprehensión intuitiva y holística del sentido global de los discursos.
    
    \item \textbf{Codificación Abierta (Primer Nivel) y Horizontalización (\textit{Horizonalization}):}
    Segmentación sistemática de los testimonios en unidades de análisis o de significado (frases, expresiones y párrafos relevantes). A cada unidad se le asigna igual valor epistémico inicial (horizontalización), eliminando enunciados repetitivos para identificar y codificar inductivamente las \textbf{categorías emergentes}.
    
    \item \textbf{Codificación Axial y Selectiva (Segundo Nivel) -- Esencia compartida vs. Divergencias individuales:}
    Agrupamiento y articulación de las categorías emergentes en temas centrales y redes semánticas complejas. En estricto cumplimiento de Hernández-Sampieri et al. \cite{sampieri2018} (p. 494), se construye la matriz dialéctica de análisis:
    \begin{itemize}
        \item \textit{Categorías esenciales o comunes (la esencia compartida):} Lo que todos o la mayoría de los usuarios tienen en común respecto al acceso y la calidad del trato humano.
        \item \textit{Categorías diferentes o divergentes (variaciones individuales):} Visiones disonantes, matices y discrepancias generacionales o de servicio asistencial.
    \end{itemize}
    
    \item \textbf{Contextualización Fenomenológica en el mundo de la vida (\textit{Lebenswelt}):}
    Elaboración integrada de dos descripciones narrativas complementarias:
    \begin{itemize}
        \item \textit{Descripción Textural:} Relato exhaustivo de \textbf{lo que} experimentaron los usuarios respecto al fenómeno.
        \item \textit{Descripción Estructural:} Comprensión profunda de \textbf{cómo} vivenciaron el fenómeno en términos de las cuatro dimensiones existenciales: \textbf{temporalidad} (tiempo vivido), \textbf{espacio} (lugar vivido), \textbf{corporalidad} (cuerpo doliente vivido) y \textbf{contexto relacional} (relación interpersonal e intersubjetiva).
    \end{itemize}
    
    \item \textbf{Síntesis narrativa de la esencia y confirmación testimonial (\textit{Member Checking}):}
    Redacción de la narrativa sintética final que transmite la esencia universal de la experiencia vivida, seguida de la devolución del reporte descriptivo a un grupo de informantes clave para corroborar que sus vivencias fueron capturadas con fidelidad suprema.
\end{enumerate}

\newpage
